\begin{proposition}[Attainment on the closed admissible family]\label{prop:attain}
For the declared closed hard predicate $\dot x\geq0$, a nonnegative feasible witness implies that the maximum primary value is attained. For any nonempty closed primary-retention superlevel, the secondary effort minimum is attained.
\end{proposition}
\begin{proof}
Consider a sequence in the closed admissible family. Compactness of $\mathcal B$ gives a convergent subsequence. Lemma~\ref{lem:regular} supplies uniform positive denominator bounds along the sequence and its limit, so every force and rate quantity converges uniformly on each normalized phase. The finitely many geometric maxima and minima are continuous; inequalities for every time are preserved under this convergence. Endpoint equalities and the ground law are continuous in $q$. Thus the admissible set is closed in a compact box. For finitely many continuous margin functions on compact normalized-time domains, $|\inf m(q,\cdot)-\inf m(q',\cdot)|\leq\sup|m(q,\cdot)-m(q',\cdot)|$, which proves continuity of $\eta$. The extreme-value theorem applies to $\eta$ and, on a closed retention set, to the continuous integral $\mathcal E$. The computed witness has strictly positive forward speed, a stronger property than the closed model constraint.
\end{proof}
